\section{Scalar compactness and critical normalization}
\label{sec:r10-foundations}
The Fourier transform is unitary, $D=-i\nabla$, and
$P(u)=\operatorname{Im}\int\overline u\nabla u=\int\xi|\widehat u|^2$.
The real Hilbert pairing is the real part of the complex $L^2$ pairing.
Throughout the local-power arguments, $N\ge1$, $m>1$, $s=m/2$,
$r=m/N$, and $p=2+2r$. A principal symbol $a$ is real, even,
$m$-homogeneous, $C^{N+2}$ off zero, and satisfies
$c|\xi|^m\le a(\xi)\le C|\xi|^m$.
Set $A(u)=\int a|\widehat u|^2$, $M(u)=\|u\|_2^2$, and
$E(u)=A(u)/2-\|u\|_p^p/p$. The sharp critical mass $M_c$ is defined by
\begin{equation}\label{eq:r10-basic-GN}
 \|u\|_p^p\le\frac p2\left(\frac{M(u)}{M_c}\right)^r A(u).
\end{equation}
The traveling-wave convention $u(t,x)=e^{i\omega t}U(x-vt)$ gives
$[a(D)+\omega+iv\cdot\nabla]U-|U|^{2r}U=0$.

\begin{lemma}[Normalized optimizer compactness]
The inequality \eqref{eq:r10-basic-GN} is attained. If
$M(f_n)\to M_c$, $A(f_n)\to1$, and $\|f_n\|_p^p\to p/2$, then
some translated subsequence converges strongly in $H^s$ to an optimizer.
The normalized optimizer set is compact modulo translations and phases.
\end{lemma}
\begin{proof}
Sharp Sobolev interpolation gives a finite positive best constant, and
amplitude and spatial scaling normalize a maximizing sequence as stated.
Bounded energy-space sequences that vanish locally in mass converge to
zero in $L^p$: use a bounded-overlap cover by unit balls, local Sobolev
interpolation, and $2<p<2N/(N-m)$ when $N>m$ (all finite $p$ otherwise).
Thus translations give a nonzero weak limit $Q$ and almost-everywhere
convergence. Write $M_1=M(Q)>0$, $A_1=A(Q)>0$, $M_2=M_c-M_1$,
and $A_2=1-A_1$. Quadratic splitting and Brezis--Lieb give
\[
 M_c^r\le M_1^r A_1+M_2^r A_2,\qquad A_1+A_2=1.
\]
If both masses are positive, the right side is strictly below $M_c^r$.
Consequently $M_1=M_c$ and then $A_1=1$. Weak convergence plus
convergence of both quadratic norms gives strong $H^s$ convergence.
The same proof applies to a sequence of optimizers. Constant phase is
proved by the density-norm lemma below; this compactness argument itself
does not need a positivity-preserving kinetic operator.
\end{proof}

